\section{Limitations}
\label{sec:limitations}
\label{sec:guarantee}
\paragraph{Coordinate semantics for automatic inference.}
TEPL supports coordinate transformations such as static slice and reverse,
but its current semantic language cannot describe every tensor access pattern.
Dynamic slice requires tensor-valued indices; pad and concatenate require
coordinate-dependent branches; and general reshape requires coordinate
flattening and unflattening using division and modulo. Without these operation
semantics, TEPL cannot automatically infer attributes or semantically verify
rewrites involving them. Extending the coordinate language and its associated
reasoning is future work.

These restrictions concern semantic verification and inference. Users can
still express structural rewrites for these operations, specify applicability
conditions, and construct RHS attributes explicitly using TEPL's attribute
expressions. Manual authoring therefore remains useful beyond the automatic
inference fragment: TEPL can match patterns, evaluate guards and attribute
programs, and build replacement graphs. Such structural rewrites do not carry
the tensor-equivalence guarantee of semantically verified rules.

\paragraph{Completeness of the model and limits of the implementation.}
Theorem~\ref{thm:complete-synthesis} establishes completeness for the exact-real,
read-once fragment in Section~\ref{sec:formal-scope}. It excludes zero and
singleton extents, fixed or shared RHS unknowns, repeated inputs within a
component, additive terms with identical input sets, constants, subtraction,
cancellation, and operations outside Table~\ref{tab:formal-operations}.
TEPL's language allows more than this fragment, including explicit attribute
derivations and descriptor reuse.
These features do not inherit its completeness guarantee.

The current implemented fragment is described in Section~\ref{sec:implementation}.
A correct rule can remain unproved because coordinate correspondence is
ambiguous, a metadata relation or coupled index equation is unsupported,
the required arithmetic laws are unavailable, or the verifier cannot
establish the guard implication. Multiple read-coordinate assignments must
agree before acceptance. Algebraic reassociation alone does not ensure
strict floating-point equivalence.
Verifying a manually supplied constructor does not guarantee that general
synthesis can discover it. Future work can expand this reasoning and the
implemented fragment without changing the theorem's stated scope.

\paragraph{Implementation soundness.}
The theorem proves constructive soundness for the mathematical algorithm;
the current compiler, symbolic translation, and broader verifier are not
formally proved sound. A full implementation proof must show that source
extraction represents the logical axes faithfully, symbolic membership and
ordering agree with the mathematical passes, and emitted programs compute
the intended attributes. It also depends on correct operation definitions,
coordinate transport, metadata and index checks, and evaluation of guards
and partial attribute programs. Reconstructing a completed rule in a fresh
context separates verification from synthesis conclusions, but both paths
share semantic infrastructure; this separation is not itself a soundness proof.

\paragraph{Reasoning cost and scalability.}
Long compositions and overlapping axis relationships can increase symbolic
partitions and normalized-expression size. Verification and inference use
finite proof budgets, so a rule may remain unresolved when reasoning exceeds
those budgets. The XLA evaluation in Section~\ref{sec:evaluation} measures
coverage; it does not establish scalability across arbitrary composition
depths or expression sizes. Future work includes systematic scaling experiments
and more efficient symbolic representations and proof search.
